%------------------------------------------------------------------------------%
\section{Integração por frações parciais.}

\begin{itemize}
  \item Seja $F$ a primitiva mais geral de $f$, ou seja,
        \(\int f(x)dx= F(x)+c\), então
        \(\int f(kx) dx= \frac{F(kx)}{k}+C\).
  \item Além disso, \(\int f(kx+b)dx= \frac{F(kx+b)}{k}+C\).
  \item \(\int \cos(5x) dx= \frac{\sin(5x)}{5}+C\).
  \item \(\int e^{8x}= \frac{e^{8x}}{8}+C\)
  \item \(\int \frac{1}{x+3}dx= \ln\abs{x+3}+C\).
\end{itemize}

Suponha que queremos calcular a integral $\int \frac{1}{x^2+a^2}dx$.
Podemos reescrever
\[
  \int \frac{1}{x^2+a^2}dx
  = \int \frac{\frac{1}{a^2}}{\left( \frac{x}{a}\right)^2+1 }dx
  = \frac{1}{a^2} \int \frac{1}{\left( \frac{x}{a}\right)^2+1} dx
\]
Usando a substituição $u= \frac{x}{a}$, temos que
$du=\frac{1}{a}dx \Longrightarrow dx=adu$.
Assim,
\[
  \int \frac{1}{x^2+a^2}dx
  = \frac{1}{a} \int \frac{du}{u^2+1}
  = \frac{1}{a} \arctan(u)+K
  =\frac{1}{a} \arctan\left(\frac{x}{a}\right)+K.
\]
Portanto,
\[
  \int \frac{1}{x^2+a^2}dx=\frac{1}{a} \arctan\left(\frac{x}{a}\right)+K.
\]
Iremos usar essa integral ao longo dos exemplos.

%------------------------------------------------------------------------------%
\begin{theory}[Método das frações parciais]
\begin{itemize}
  \item Método para integrar funções racionais, ou seja,
        $f(x)=\frac{P(x)}{Q(x)}$, onde $P$ e $Q$ são polinômios.

  \item Se o grau de $P$ for maior ou igual que o grau de $Q$,
        precisamos de uma etapa preliminar que é efetuar a divisão
        de $P$ por $Q$.

  \item Feito isso, o próximo passo é fatorar $Q(x)$ ao máximo;

  \item Expandir $f$ em frações parciais e determinar os coeficientes.
\end{itemize}
\end{theory}
%------------------------------------------------------------------------------%

Casos a serem analisados.

%------------------------------------------------------------------------------%
\begin{theory}[Caso 1]
  $Q(x)$ é o produto de $m$ fatores lineares distintos, isto é,
  $Q(x)=(a_1x+b_1)(a_2x+b_2)\dots(a_mx+b_m)$.
  \[
    \frac{P(x)}{Q(x)}
    = \frac{A_1}{(a_1x+b_1)}
    + \frac{A_2}{(a_2x+b_2)}
    + \cdots
    + \frac{A_m}{(a_mx+b_m)}
  \]
\end{theory}
%------------------------------------------------------------------------------%

Por exemplo,
\[
  \frac{x-1}{x(x-1)(x+1)}= \frac{A}{x}+ \frac{B}{x-1}+ \frac{C}{x+1}
\]

%------------------------------------------------------------------------------%
\begin{theory}[Caso 2]
  $Q(x)$ é o produto de $m$ fatores lineares e alguns dos
  fatores são repetidos. Para simplificar, suponha que o primeiro
  fator linear $(a_1x+b_1)$ seja repetido
  \[
    \frac{A_1}{(a_1x+b_1)}
    + \frac{A_2}{(a_1x+b_1)^2}
    + \cdots
    + \frac{A_i}{(a_1x+b_1)^i}
  \]
\end{theory}
%------------------------------------------------------------------------------%

Por exemplo,
\[
  \frac{x^3-x-2}{x^3(x-1)}
  = \frac{A}{x}
  +\ frac{B}{x^2}
  + \frac{C}{x^3}
  + \frac{D}{x-1}
\]

%------------------------------------------------------------------------------%
\begin{theory}[Caso 3]
  $Q(x)$ contém fatores quadráticos irredutíveis que não se repetem.
  Se $Q(x)$ contém um fator $ax^2+bx+c$, onde $b^2-4ac<0$,
  então $\frac{P(x)}{Q(x)}$ tem um fator na forma
  \[
    \frac{Ax+B}{ax^2+bx+c}
  \]
\end{theory}
%------------------------------------------------------------------------------%

Por exemplo,
\[
  \frac{x}{(x-2)(x^2+1)(x^2+4)}
  = \frac{A}{x-2}
  + \frac{Bx+C}{x^2+1}
  + \frac{Dx+E}{x^2+4}
\]

%------------------------------------------------------------------------------%
\begin{theory}[Caso 4]
  $Q(x)$ contém fatores quadráticos irredutíveis repetidos.
  Se $Q(x)$ contém um fator $(ax^2+bx+c)^r$, onde $b^2-4ac<0$,
  então ao invés de uma fração parcial temos a seguinte decomposição
  \[
    \frac{A_1x+B_1}{ax^2+bx+c}
    + \frac{A_2x+B_2}{(ax^2+bx+c)^2}
    + \cdots
    + \frac{A_rx+B_r}{(ax^2+bx+c)^r}
  \]
\end{theory}
%------------------------------------------------------------------------------%

Por exemplo,
\[
  \frac{2x}{(x-9)(x^2+4)^2}
  = \frac{A}{x-9}
  + \frac{Bx+C}{x^2+4}
  + \frac{Dx+E}{(x^2+4)^2}
\]

%------------------------------------------------------------------------------%
\begin{example}
  Calcule a integral \(\int \frac{x^3+x}{x-1}\)
  \solution
  Uma vez que o grau do numerador é maior que o grau do denominador, precisamos
  fazer a etapa preliminar, que é efetuar a divisão de polinômios.
  \begin{align*}
    \int \frac{x^3+x}{x-1}
    & = \int \left( x^2+x+2 + \frac{2}{x-1} \right) dx \\
    & = \frac{x^3}{3}
      + \frac{x^2}{2}
      + 2x
      + 2 \ln\abs{x-1}
      + C
  \end{align*}
\end{example}
%------------------------------------------------------------------------------%

Nesse exemplo a etapa preliminar já resolveu o nosso problema.

%------------------------------------------------------------------------------%
\begin{example}
  Calcule \(\int \frac{x^2+2x-1}{2x^3+3x^2-2x}dx\)
  \solution
  Nesse caso, o grau do numerador é menor que o grau do denominador, então não
  precismos fazer a etapa preliminar. Então, o próximo passo é fatorar o
  denominador.
  \[
    2x^3+3x^2-2x
    = x(2x^2+3x-2)
    = 2x(x- \frac{1}{2})(x+2)
    = x(2x-1)(x+2)
  \]
  Temos o primeiro caso, uma vez que o denominador de produto de fatores
  lineares distintos. Assim, a decomposição em frações parciais do integrando
  tem a forma.
  \[
    \frac{x^2+2x-1}{2x^3+3x^2-2x}
    = \frac{A}{x}
    + \frac{B}{2x-1}
    + \frac{C}{x+2}
  \]
  A fim de determinamos os valores de $A,B$ e $C$, multiplicamos ambos os
  lados da equação pelo MMC dos denominadores, $x(2x-1)(x+2)$, obtendo
  \begin{align*}
    x^2+2x-1
    & = A(2x-1)(x+2)+Bx(x+2)+ Cx(2x-1) \\
    & = (2A+B +2C)x^2+ (3A+2B-C)x-2A
  \end{align*}
  Usando a identidade de polinômios, igualamos os coeficientes de cada termo,
  por exemplo, o coeficiente de $x^2 $ do lado direito, $2A+B +2C$ deve ser
  igual ao coeficiente de $x^2$ dolado esquerdo, ou seja, $1$. Prosseguindo
  assim, temos o seguinte sistemas de equação
  \begin{align*}
    2A+B +2C & = 1 \\
    3A+2B-C  & = 2 \\
    -2A      & =-1
  \end{align*}
  Resolvendo, encontramos
  $A= \frac{1}{2}$,
  $B= \frac{1}{5}$ e
  $C=-\frac{1}{10}$,
  e assim
  \begin{align*}
    \int \frac{x^2+2x-1}{2x^3+3x^2-2x} dx
    & = \int \left(
                \frac{1}{2}  \frac{1}{x}
              + \frac{1}{5}  \frac{1}{2x-1}
              - \frac{1}{10} \frac{1}{x+2}
             \right) dx \\
    & = \frac{1}{2}  \ln\abs{x}
      + \frac{1}{10} \ln\abs{2x-1}
      - \frac{1}{10} \ln\abs{x+2}
      + C
  \end{align*}
\end{example}
%------------------------------------------------------------------------------%

%------------------------------------------------------------------------------%
\begin{example}
  Calcule \(\int \frac{4x}{x^3-x^2-x+1}dx\).
  \solution
  Note que novamente não precisamos fazer a etapa preliminar.
  As raízes do polinômio $x^3-x^2-x+1$ são $x=1$,
  com multiplicidade 2 e $x=-1$. Assim,
  \[
    x^3-x^2-x-2=(x-1)^2(x+1)
  \]
  Como o fator linear $x-1$ aparece duas vezes, temos que a decomposição
  em frações parciais é dada por
  \[
    \frac{4x}{x^3-x^2-x-2}
    = \frac{A}{x-1}
    + \frac{B}{(x-1)^2}
    + \frac{C}{x+1}
  \]
  Multiplicando essa expressão pelo mínimo múltiplo comum $(x-1)^2(x+1)$, obtemos
  \begin{align}
    4x & = A(x-1)(x+1)+B(x+1)+C(x-1)^2 \\
       & = (A+C)x^2+(B-2C)x+(-A+B+C)
  \end{align}
  Assim, para acharmos os coeficientes $A,B$ e $C$ precisamos resolver o sistema
  \begin{align*}
    A+C    & = 0 \\
    B-2C   & = 4 \\
    -A+B+C & = 0
  \end{align*}
  Cuja solução é $A=1$, $B=2$ e $C=-1$. Assim,
  \begin{align*}
    \int \frac{4x}{x^3-x^2-x-2}dx
    & = \int \left[ \frac{1}{x-1}+ \frac{2}{(x-1)^2}-\frac{1}{x+1} \right] dx \\
    & = \ln \abs{x-1}- \frac{2}{x-1}- \ln\abs{x+1}+C \\
    & = \ln \abs{\frac{x-1}{x+1}} - \frac{2}{x-1}+C
  \end{align*}
\end{example}
%------------------------------------------------------------------------------%

%------------------------------------------------------------------------------%
\begin{example}
  Encontre \(\int \frac{2x^2-x+4}{x^3+4x}dx\).
  \solution
  Note que $x^3+4x=x(x^2+4)$, então temos que o denominador é o produto de um
  fator linear com um fator quadrático irredutível. Portanto,
  \[
    \frac{2x^2-x+4}{x^3+4x}
    = \frac{2x^2-x+4}{x(x^2+4)}
    = \frac{A}{x}+ \frac{Bx+C}{x^2+4}
  \]
  Multiplicando por $x(x^2+4)$, temos
  \begin{align*}
    2x^2-x+4
    & = A(x^2+4)+(Bx+C)x \\
    & = (A+B)x^2+Cx+4A
  \end{align*}
  Igualando os coeficientes, temos o seguinte sistema
  \[
    A+B=2 \qquad C=-1 \qquad 4A=4
  \]
  Assim, temos que $A=1$, $B=1$ e $C=-1$. Portanto,
  \begin{align*}
    \int \frac{2x^2-x+4}{x^3+4x}dx
    & = \int \left( \frac{1}{x}+ \frac{x-1}{x^2+4} \right)dx \\
    & = \int \frac{1}{x}dx + \int \frac{x}{x^2+4} dx - \int \frac{1}{x^2+4} dx \\
    & = \ln\abs{x}
      + \frac{1}{2} \ln(x^2+4)
      - \frac{1}{2} \arctan\left(\frac{x}{2}\right)
      + C
  \end{align*}
  Onde para resolver a integral \(\int \frac{x}{x^2+4} dx\) usamos a integração
  por substituição simples, tomando $u=x^2+4$, de modo que $du=2xdx$ e para
  calcular \(\int \frac{1}{x^2+4}\) usamos a fórmula
  \[
    \int \frac{1}{x^2+a^2}dx
    = \frac{1}{a} \arctan\left(\frac{x}{a}\right)+K,
  \]
  com $a=4$.
\end{example}
%------------------------------------------------------------------------------%

%------------------------------------------------------------------------------%
\begin{example}
  Calcule \(\int \frac{\sin(x)}{\cos^2(x)-3\cos(x)}dx\).
  \solution
  Note que aqui não temos,a princípio uma função racional. Vamos começar a
  resolver essa integral usando a substituição simples $u=\cos(x)$, de modo que
  $du=-\sin(x) dx \to -du=\sin(x)dx$. Assim,
  \[
    \int \frac{\sin(x)}{\cos^2(x)-3\cos(x)}dx
    =  \int \frac{-du}{u^2-3u}
    = -\int \frac{1}{u(u-3)}du
  \]
  Agora temos uma função racional em $u$. Expandindo em frações parciais, temos
  \[
    \frac{1}{u(u-3)}
    = \frac{A}{u}+ \frac{B}{u-3}
  \]
  Multiplicando pelo mínimo múltiplo comum $u(u-3)$ temos:
  \[
    1
    = A(u-3)+Bu
    = (A+B)u-3A
  \]
  Assim, temos que $A=-\frac{1}{3}$ e $B=\frac{1}{3}$. Portanto,
  \begin{align*}
    \int \frac{\sin(x)}{\cos^2(x)-3\cos(x)}dx
    & =  \int \frac{-du}{u^2-3u} \\
    & = -\int \frac{1}{u(u-3)}du \\
    & = -\int -\frac{1}{3} \frac{1}{u}+\frac{1}{3} \frac{1}{u-3} \\
    & = \frac{1}{3} \ln\abs{u}
      - \frac{1}{3} \ln\abs{u-3} + C \\
    & = \frac{1}{3} \ln \abs{\frac{u}{u-3}} +C \\
    & = \frac{1}{3} \ln \abs{\frac{\cos(x)}{\cos(x)-3}} +C
  \end{align*}
\end{example}
%------------------------------------------------------------------------------%
